\chapter{Introduction}\label{chap:1}
\label{ch:intro}

\section{What REGEN is}
\label{sec:ch01-what}

REGEN (Response Engineering with GEnetic algorithms and Neural networks) is a Windows desktop application for the
design and analysis of experiments. It brings together, in one project file, the classical tools of design of
experiments (DOE) and response surface methodology (RSM) \citep{montgomery2017,myers2016}, artificial neural
networks (ANN) as an alternative or complementary response model, and multi-objective optimization with the
evolutionary algorithms NSGA-II and NSGA-III \citep{deb2002,deb2014}.

REGEN is meant for engineers, researchers and students who plan experiments, collect data in the laboratory or the
plant, and want to turn the data into a model and an optimum that can be confirmed. Typical questions are:
\begin{itemize}
  \item Which design and how many runs do I need to screen four factors, or to fit a quadratic model in three
        factors?
  \item Are my data good enough to build a model, and which runs should I add?
  \item Which factors and interactions are significant, and how well does the model predict new runs?
  \item Is an ANN better than the polynomial model for my data, and do I have enough data for an ANN?
  \item Which settings give the best compromise between conflicting responses, and do the RSM and ANN models agree
        on that compromise?
  \item How sensitive is the optimum to each factor, and does a confirmation run fall inside the prediction
        interval?
\end{itemize}

The program runs entirely on the local computer. No internet connection, no account and no Python installation are
needed. All results can be exported (Excel, CSV, PDF, clipboard), so every number and plot can be reproduced or
redrawn outside REGEN.

\section{The workflow}
\label{sec:ch01-workflow}

The pages of REGEN follow the order in which an experimental study is usually carried out. The navigation tree on
the left of the main window groups them by stage (Project, Design, Analysis, Optimization, Tools). A complete study
passes through the following steps; each step is described in its own chapter.

\begin{enumerate}
  \item \textbf{Design.} Choose a design with \ui{New Design} or let the \ui{Design Guide} recommend a design and a
        number of runs for the goal (screening, characterization, RSM optimization, ANN modeling, or RSM and ANN
        together). Evaluate the design before any run is made (\ui{Design Evaluation}: degrees of freedom, power,
        aliasing, variance inflation, fraction of design space). See Chapter~\ref{chap:3}.
  \item \textbf{Data.} Run the experiments in the randomized run order and type or paste the responses in the
        \ui{Data} page, or import existing (historical) data from Excel or CSV. Runs can be excluded from the
        analysis without deleting them, and the design can be augmented. See Chapter~\ref{chap:4}.
  \item \textbf{Readiness.} The \ui{Data Readiness} page scores the data for modeling and for optimization and
        proposes extra runs (D-optimal, space-filling, replicates, confirmation runs) that can be added to the data
        sheet with one click. See Chapter~\ref{chap:4}.
  \item \textbf{RSM.} Fit polynomial, Scheff\'e mixture, logistic/Poisson or split-plot models on the
        \ui{RSM Analysis} page: fit summary, Type~III ANOVA, equations, diagnostics, Box-Cox and factorial effects.
        See Chapter~\ref{chap:5}.
  \item \textbf{ANN.} Assess data adequacy for an ANN, search the network architecture automatically by
        cross-validation, train the final network and check with a learning curve whether more data would help.
        RSM or ANN can be selected separately for each response. See Chapter~\ref{chap:7}.
  \item \textbf{Graphs.} Contour, 3D surface, one-factor, interaction, cube, perturbation, ternary, trace,
        overlay, desirability and propagation-of-error plots of the selected model (\ui{Model Graphs}). See
        Chapter~\ref{chap:6}.
  \item \textbf{Optimization.} Desirability optimization of one or more responses, or the Pareto front of two or
        more objectives with NSGA-II/NSGA-III, with an automatic stopping rule, a population-size check and a
        comparison of the fronts obtained with RSM and with ANN. See Chapter~\ref{chap:8}.
  \item \textbf{Sensitivity.} Parametric studies, global sensitivity (standardized regression coefficients or
        Sobol indices), local sensitivity (tornado chart) and the trade-off along the Pareto front. See Chapter~\ref{chap:9}.
  \item \textbf{Prediction and confirmation.} Predict the responses at the chosen settings with confidence and
        prediction intervals, run confirmation experiments and enter their mean to see whether the model is
        confirmed. See Chapter~\ref{chap:10}.
\end{enumerate}

Chapter~\ref{chap:11} follows one case study through the whole workflow.

\section{The four diagnostic checks}
\label{sec:ch01-checks}

When an ANN is used next to (or instead of) RSM, and when the optimum is searched with an evolutionary algorithm,
several decisions are usually made by habit: the size of the network, whether the data are sufficient, when the
algorithm has run long enough, and whether the optimum depends on the chosen model. REGEN turns these decisions
into four explicit checks with numerical criteria. Each check answers one question, applies a stated pass
criterion and tells the user what to do when the criterion is not met.
Table~\ref{tab:ch01-checks} summarizes them; the criteria are those implemented in the program.

{\footnotesize
\begin{longtable}{@{}p{0.11\linewidth}p{0.16\linewidth}p{0.36\linewidth}p{0.22\linewidth}p{0.03\linewidth}@{}}
\caption{The four diagnostic checks of REGEN. $E$ denotes the normalized mean squared error (MSE divided by the
variance of the response; $R^2 = 1-E$), $E_{te}$ the cross-validated test error, $n^*$ the estimated number of
training runs, $n_{full}$ the largest training subset of the learning curve, HV the hypervolume and $w$ the
generation window.}\label{tab:ch01-checks}\\
\toprule
Check & Question & Pass criterion & Action when it fails & Ch. \\
\midrule
\endfirsthead
\toprule
Check & Question & Pass criterion & Action when it fails & Ch. \\
\midrule
\endhead
\bottomrule
\endfoot
A. Architecture search &
Which network (layers and neurons) generalizes best for the available data? &
Networks are grown one neuron count at a time while the cross-validated test error improves by at least 1\,\%
(stop after 2 failures in a row; a new layer is kept only if it improves the best test error by at least 1\,\%).
Each candidate is classified Good fit, Overfitting or Underfitting. The selected network is the smallest (fewest
weights) Good candidate whose test error is within $\max(\mathrm{SE},\,0.002)$ of the lowest test error
(one-standard-error rule). &
If no candidate is a Good fit, all candidates are used and the status is reported; if the ANN test $R^2$ is not
above the RSM value on the same folds, the program states that RSM predicts new data as well as or better than
ANN: keep RSM, or add runs. & 7 \\
\addlinespace
B. Learning curve &
Would the ANN still improve with more runs? &
A power law $e(n)=a+b\,n^{-c}$ is fitted to the median test error; $n^*$ is the training size at which the test
$R^2$ is within 0.01 of its limit. ``The data is sufficient'' if the median $n^*\le n_{full}$; ``more data would
still help'' if $n_{full}<n^*\le5\,n_{full}$; ``far from sufficient'' if $n^*>5\,n_{full}$. Fewer than 25 runs:
``too little data''; fewer than 12 runs: not computed. &
Add about the estimated number of runs (gradually, for example half first) using the space-filling suggestions of
Data Readiness, then repeat; with far too little data use RSM, fewer factors or narrower ranges. & 7 \\
\addlinespace
C. Hypervolume stopping rule &
Has NSGA-II/III run enough generations, and is the population large enough? &
Generations: the mean hypervolume of the last $w=20$ generations changes by less than
$\max(\mathrm{tol},\min(\mathrm{noise},1\,\%))$ (tol = 0.1\,\%) compared with the previous $w$ generations, and
the ideal point moves by less than 1\,\% of the range; the run stops when this has held for $w$ more generations
(at least 30 generations). Population: two other seeds change the HV by less than 1\,\% and a doubled population
adds less than 1\,\%. &
Increase the maximum number of generations; increase the population (for example to $2N$) and check again; for
four or more objectives use NSGA-III. & 8 \\
\addlinespace
D. ANN vs RSM fronts &
Does the optimum depend on the choice of model? &
Both fronts are re-predicted with the other model. $\Delta$ = largest (over objectives) mean absolute difference
between the RSM and ANN predictions on the front solutions, in \% of the response data range; $dx$ = largest
distance between the two TOPSIS choices in \% of the factor range. Agree: $\Delta<5\,\%$ and $dx<15\,\%$. Partial
agreement: $\Delta<10\,\%$. Otherwise the models differ. &
Do not rely on one model: make confirmation runs at both TOPSIS points, check the region of the largest
differences, and trust the model with the better predictive accuracy (RSM predicted $R^2$, ANN cross-validated
$Q^2$). & 8 \\
\end{longtable}}

The checks do not replace engineering judgement. They make the reasons for a decision visible and reproducible, and
they are reported in the program and in the PDF report so that a reader of the results can see on which basis a
network, a stopping point or an optimum was accepted.

\section{Conventions used in this manual}
\label{sec:ch01-conventions}

\paragraph{User-interface labels.} Names of pages, buttons, options and menu items are printed in a sans-serif font,
for example \ui{Run Optimization} or \ui{RSM Analysis}, exactly as they appear on the screen. Menu paths are written
as \menupath{File $\rightarrow$ Case Studies\dots}. A label that ends in ``\dots'' opens a dialog.

\paragraph{Factors, responses and letters.} Factors are named with capital letters A, B, C, \dots{} in the order in
which they are defined; the letter I is skipped to avoid confusion with the number 1, so the letters run A--H, J--Z
(at most 25 factors). Responses are numbered R1, R2, \dots{} and shown as, for example, \ui{R1: Yield (\%)}.

\paragraph{Actual and coded units.} Factor values are entered and displayed in \emph{actual units} (for example
\textdegree C or minutes). Internally, and in the coded model equation, a numeric factor $x$ with low level $L$ and
high level $H$ is used in \emph{coded units}
\begin{equation}
  z = \frac{x-(L+H)/2}{(H-L)/2}, \qquad -1 \le z \le +1 \text{ inside the design range},
  \label{eq:ch01-coded}
\end{equation}
so that the low level is $-1$, the center is 0 and the high level is $+1$. Coefficients in coded units can be
compared with each other directly; coefficients in actual units cannot, because they depend on the units.
Categoric factors are stored as level indices and enter the models with effect coding (a factor with $L$ levels
gives $L-1$ columns).

\paragraph{Mixture components and pseudo-components.} In a mixture the components $x_i$ add up to a fixed total $T$
(for example 100\,\%), and each component has a lower limit $L_i$. REGEN works in \emph{pseudo-components}
\begin{equation}
  z_i = \frac{x_i - L_i}{T - \sum_j L_j}, \qquad \sum_i z_i = 1,
  \label{eq:ch01-pseudo}
\end{equation}
which map the constrained mixture region onto a simplex \citep{cornell2002,piepel1982}. Scheff\'e mixture
models \citep{scheffe1958} and ternary plots are expressed in pseudo-components; the data sheet shows actual amounts
unless \ui{Show pseudo-components} is ticked.

\paragraph{Numbers.} Both the decimal point and the decimal comma are accepted wherever numbers are typed or pasted.
Copied tables use the decimal separator of the Windows regional settings (a second menu item copies with the other
separator). In this manual the decimal point is used.

\paragraph{Case studies.} Most examples use the twelve built-in case studies (\menupath{File $\rightarrow$ Case
Studies\dots}, Section~\ref{sec:ch02-cases}). Their data are generated with fixed random seeds, so the numbers shown
in this manual can be reproduced exactly. The screenshots were made with one CPU core and the default seeds.

\section{How this manual is organized}
\label{sec:ch01-organization}

\begin{itemize}
  \item Chapter~\ref{chap:2}, \emph{Getting Started}: installation, the main window, menus, project files, case studies,
        import and export.
  \item Chapter~\ref{chap:3}, \emph{Designs}: all design types, the Design Guide and Design Evaluation.
  \item Chapter~\ref{chap:4}, \emph{Data}: the data sheet, augmentation, simulation, Data Readiness and normalization.
  \item Chapter~\ref{chap:5}, \emph{RSM Analysis}: model fitting, ANOVA, diagnostics, transformations, GLM and split-plot
        analysis.
  \item Chapter~\ref{chap:6}, \emph{Model Graphs}.
  \item Chapter~\ref{chap:7}, \emph{Artificial Neural Networks}, including checks A and B.
  \item Chapter~\ref{chap:8}, \emph{Optimization}: desirability, NSGA-II/III, checks C and D.
  \item Chapter~\ref{chap:9}, \emph{Sensitivity and Parametric Study}.
  \item Chapter~\ref{chap:10}, \emph{Prediction and Confirmation}.
  \item Chapter~\ref{chap:11}, \emph{worked examples} from design to confirmation, including a multi-objective study
        of a polygeneration plant.
  \item Chapter~\ref{chap:12}, \emph{Troubleshooting and Frequently Asked Questions}.
  \item Appendix~A, \emph{Reference}: keyboard shortcuts, the project file format, constants and thresholds, and a
        glossary.
\end{itemize}

Each feature section follows the same pattern: purpose, how to use it, options, output and how to read it, the
theory as implemented (with the constants used by the program), an example and a note on data export.

\section{How to cite REGEN}
\label{sec:ch01-cite}

REGEN is developed by Erdiyanto Munandar and Nasruddin, Department of Mechanical Engineering, Faculty of
Engineering, Universitas Indonesia. The source code is available at \url{https://github.com/chumporium/REGEN} under
the MIT License. Every released version is archived in Zenodo:
\begin{itemize}
  \item version 7.0.0: DOI \href{https://doi.org/10.5281/zenodo.22944177}{10.5281/zenodo.22944177};
  \item all versions (concept DOI, always resolves to the latest version):
        \href{https://doi.org/10.5281/zenodo.22944176}{10.5281/zenodo.22944176}.
\end{itemize}
Please cite the version that was used for the work, for example:

\begin{quote}
Munandar, E., and Nasruddin. 2026. REGEN: Response Engineering with GEnetic algorithms and Neural networks, version
7.0.0. Zenodo. \url{https://doi.org/10.5281/zenodo.22944177}
\end{quote}

A BibTeX entry:
\begin{verbatim}
@software{regen700,
  author    = {Munandar, Erdiyanto and Nasruddin},
  title     = {{REGEN}: Response Engineering with GEnetic algorithms
               and Neural networks},
  version   = {7.0.0},
  year      = {2026},
  publisher = {Zenodo},
  doi       = {10.5281/zenodo.22944177},
  url       = {https://doi.org/10.5281/zenodo.22944177}
}
\end{verbatim}
The repository also contains the citation metadata (file \texttt{CITATION.cff}, shown by GitHub as \emph{Cite this
repository}). When a method implemented in REGEN is essential for the results (for example the desirability
function, NSGA-II or Sobol indices), please cite the original method as well; the references are given in the
corresponding chapters.

REGEN uses open-source libraries that keep their own licenses; the user interface uses PySide6 (Qt for Python),
which is licensed under the LGPL-3.0. Versions 1 to 6 of the program were developed under the name DOE Studio;
project files of those versions can still be opened.
